\section{Metrics and Statistics}\label{sec:metrics}
Per system: \emph{detection} = P(\DIV{} $\mid$ divergent label); \emph{false assurance} = P(\NDO{} $\mid$ divergent label); \emph{abstention} = P(\INS{} or failed run) over the primary denominator; \emph{false alarm} = P(\DIV{} $\mid$ no-difference label); \emph{risk--coverage} with coverage = decided/eligible and risk = errors/decided, where errors are \NDO{} on a divergent label or \DIV{} on a no-difference label; CPU per pair. Detection and false assurance are the primary family; the rest are secondary. Failed or missing runs count as abstention (intention-to-treat). Baseline outputs map DIVERGE to \DIV{} and NO\_DIFF to \NDO{}. Labels are generator labels or test-oracle truth and are not adjudicated, so ``false assurance'' here means ``no difference reported on a pair labelled divergent'', not a verified miss.

Inference unit: project. Rates are ratios with a project-cluster percentile bootstrap (95\%, $B=10000$, seed 42). Paired comparisons use a project-level sign-flip test, an exact McNemar test at pair level as a sensitivity check, and Holm correction within the family of the comparisons in this run. The +10pp rule requires a Holm-adjusted significant result and a lower confidence bound above +10pp. With few independent projects, percentile bootstrap intervals undercover; here one cluster of 29 divergent pairs (QuixBugs) dominates, so the intervals are not trustworthy. Power is simulated on the observed design (assumptions A1--A4 are in the source file) and post-hoc power at the observed effect is not reported. Table~\ref{tab:notevaluated} lists what is not evaluated.
\begin{table*}[t]
\centering\scriptsize
\caption{NOT EVALUATED register (from \path{research/evaluation/EXPLORATORY_RESULTS.md} and \path{research/evaluation/K1_K2_FEASIBILITY_REPORT.md}).}
\label{tab:notevaluated}
\begin{tabular}{p{0.27\textwidth}p{0.14\textwidth}p{0.5\textwidth}}
\hline
Item & Status & Reason \\
\hline
Confirmatory family (P1/P2; H2a, H2b, H4) & NOT EVALUATED & no sealed holdout; data non-pristine; K5 not met; no independent real projects \\
Superiority over a like-for-like comparator; K3/K4 gates; matched-coverage false assurance & NOT EVALUATED & only raw counts; no matched-coverage computation \\
RQ4 generalisation & NOT ASSESSABLE & all data development-tuned; no independent project-level run \\
RQ5 overhead & LOWER BOUND ONLY & V7-B CPU is child-process only; parent overhead and memory not recorded \\
Truth for BugsInPy units, two QuixBugs pairs, six suspect labels & NOT EVALUATED & no defensible label; never relabelled \\
ICC and design effect & UNRELIABLE & one large cluster dominates; estimate is zero \\
K1 and K2 (isolation-based) & NOT MET / NOT ASSESSABLE & fewer runnable real pairs than required; no project-level runs; no Linux isolation \\
Published tools (Shadow, Testora, Mokav, DiffTestGen, Cleverest, DiffKemp, AFL-class fuzzers) & NOT IMPLEMENTED & no comparison; no claim relative to them \\
S3, explorer, module E, Linux container isolation, declared renames, generator entries, Merkle and Supersession records & NOT IMPLEMENTED & see implementation status; they cause most INSUFFICIENT outcomes \\
Component ablation of V7-B (replay, control alternation, thresholds) & NOT EVALUATED & never run; only B1 versus B1+rep is available as a proxy \\
Sensitivity to thresholds (0.80 branch floor, repeat counts, budgets) & NOT EVALUATED & defaults are unvalidated and were not varied \\
Cross-language, cross-platform & NOT SUPPORTED & Python on macOS only \\
\hline
\end{tabular}
\end{table*}

